% REVISED: canonical-target reframe 2026-06-04
% C-COMPRESS: appendix reduced to opportunity + timing validation details
% (6 threat-family subsections, 3 anchor tables). Permutation, matched,
% leakage, audit-map, timing-sets, timing-leakage, clustered-RI, and
% direct-defendant-timing tables moved to Online Supplement C; headline
% numbers retained in prose.
\section{Validation Audits}
\label{app:validation_audits_submission}

The validation design addresses three mechanical channels.
Frequent participants have more chances to appear near any
defendant, so raw co-bidding rates are not enough; firms may
operate in the same procurement environments as CADE anchors,
giving them more opportunities to meet those anchors even under
ordinary competition; and item-level labels can mechanically
reuse the same participation history that constructs the score.
This appendix records the opportunity and timing audit logic and
the interpretation each audit licenses; the main text reports the
headline results.

Throughout, ``CADE-defendant contact'' denotes an
\emph{adjudication-anchored exposure} label: a firm shares a
tender-item with a directly adjudicated CADE defendant. It does
not assert that the firm is a cartel member, party to any
agreement, or liable for anything. The audits validate forensic
priority, not membership, damages, or proof. Online Supplement C
reports the full grids behind every audit summarized here.

\subsection{Opportunity-Cell Construction and Expected Contact}
\label{app:opp_cells_submission}
\label{app:market_opportunity_submission}
\label{app:expected_contact_submission}

\textit{Purpose.} A firm can appear near a CADE anchor simply
because it operates in environments where defendants are active;
the opportunity cells discipline this exposure before any score
is credited.

\textit{Method.} We partition the award layer into
procurement-opportunity cells $g$ and benchmark each firm against
others facing the same cells, using three nested definitions:
\textbf{COARSE} $=$ item group (two-digit item-code prefix)
$\times$ year; \textbf{MEDIUM} $=$ item group $\times$ year
$\times$ buyer/procuring-unit (the reference definition for the
headline within-stratum results); and \textbf{STRICT} $=$ item
code $\times$ year $\times$ buyer. Modality (Convite vs.\
Preg\~ao) is intentionally dropped: it is not present in the
award layer used here, and conditioning on it would further
fragment already-sparse cells (revisited as a sensitivity below).
For firm $i$, with $n_{ig}$ the tender-items held in cell $g$ and
$p_g$ the cell defendant-contact rate, \emph{expected} contact
under ordinary opportunity is $E_i = \sum_g n_{ig}\, p_g$, the
\emph{observed} contact $O_i$ is the firm's actual count of
defendant-sharing tender-items, and the \emph{excess} is
$X_i = O_i - E_i$ with standardized form
$Z_i = (O_i - E_i)/\sqrt{\sum_g n_{ig}\, p_g(1-p_g)}$. To prevent
a firm from being benchmarked against itself, $p_g$ is computed
leave-one-out.

\textit{Result.} The mean cell defendant-contact rate $\bar p_g$
is around $0.017$--$0.019$ across the three definitions; finer
cells gain resolution at the cost of support (under STRICT,
$99.5\%$ of cells hold fewer than five tender-items and $75.7\%$
are singletons). Excess concentrates where the score is high
(mean $X_i = +0.068$ at participation tier $T\ge 50$, versus
$\approx -0.006$ to $-0.000$ at lower tiers), but most of the
level of $O_i$ is explained by $E_i$: exposure does the heavy
lifting and excess is the residual the score must rank.
Common-support retention is honest about sparsity --- under
MEDIUM cells the strict support rule retains about $46\%$ of
firms (the within-stratum results use MEDIUM), and under STRICT
retention collapses to about $9\%$. The cell-construction,
expected-contact, and retention grids are in Online Supplement~C.

\subsection{Opportunity-Adjusted Validation}
\label{app:control_function_submission}

\textit{Purpose.} The control-function audit asks whether the
screening score still ranks defendant contact once opportunity
exposure $E_i$ is in the model.

\textit{Method.} We estimate a ladder of logits and rank-tests
(S0--S6) on the always-loser firms, reported in
Table~\ref{tab:control_function_submission}: raw baselines (S0,
including an exposure-only model), flexible participation (S1),
score-plus-$\log(1{+}E)$ at each cell granularity (S2), breadth
controls (S3), opportunity- and participation-fixed effects (S4),
the binary excess target $\mathbf{1}[X_i>0]$ (S5), and the
strict-support subsample (S6).

\begin{table}[htbp]
\centering
\caption{Control-Function Validation (Selected Specifications)}
\label{tab:control_function_submission}
\scriptsize
\begin{adjustbox}{max width=\textwidth,center}
\begin{threeparttable}
\begin{tabular}{lL{4.0cm}rrrr}
\toprule
Spec & Model & $N$ & Pos. & ROC AUC & PR-AUC \\
\midrule
S0 & raw score (log\,$T$) & $16{,}843$ & $651$ & $0.761$ & $0.143$ \\
S0 & raw FL14 & $16{,}843$ & $651$ & $0.688$ & $0.097$ \\
S0 & exposure-only (all AL) & $16{,}843$ & $651$ & $0.905$ & $0.300$ \\
S2 & score $+\log(1{+}E)$ [MEDIUM] & $16{,}843$ & $651$ & $0.985$ & $0.708$ \\
S2 & FL14 $+\log(1{+}E)$ [MEDIUM] & $16{,}843$ & $651$ & $0.985$ & $0.706$ \\
S2 & exposure-only logit (exposed) & $6{,}040$ & $651$ & $0.713$ & $0.300$ \\
S3 & score $+\,E\,+$ breadth & $16{,}843$ & $651$ & $0.983$ & $0.711$ \\
S4 & score $+$ opp-FE $+$ part-FE (exposed) & $6{,}040$ & $651$ & $0.707$ & $0.255$ \\
S5 & score ranks $\mathbf{1}[X_i{>}0]$ (MEDIUM) & $16{,}843$ & $475$ & $0.764$ & $0.095$ \\
S5 & FL14 ranks $\mathbf{1}[X_i{>}0]$ (MEDIUM) & $16{,}843$ & $475$ & $0.687$ & $0.073$ \\
S6 & score $+E$, strict support & $7{,}683$ & $446$ & $0.969$ & $0.671$ \\
NESTED & exposure-only logit & $6{,}040$ & $651$ & $0.713$ & --- \\
NESTED & exposure $+$ score & $6{,}040$ & $651$ & $0.723$ & --- \\
WITHIN & score $\mid$ opportunity stratum & $6{,}040$ & $651$ & $0.471$ & --- \\
WITHIN & FL14 $\mid$ opportunity stratum & $6{,}040$ & $651$ & $0.507$ & --- \\
\bottomrule
\end{tabular}
\begin{tablenotes}
\scriptsize
\item \textit{Notes:} Outcome is adjudication-anchored cobidder
contact (S5 uses the binary excess target $\mathbf{1}[X_i>0]$).
\textbf{The combined exposure-plus-score logits in S2--S3 reach
AUC $\approx 0.98$, but this is mechanically inflated and is NOT
evidence of score signal:} the exposure term $E_i$ is built from
the same defendant-contact rates that define the label, so a
model containing $E_i$ partly predicts the label by construction.
The score's \emph{isolated} contribution is not read off these
combined AUCs. Holding exposure fixed, the residual ordering is
marginal at best: the nested DeLong increment is only $+0.010$
over the exposure-only model ($p=0.013$), and the
within-opportunity-stratum AUC for the score is $0.471$ ($\approx$
chance; FL14 $0.507$). The companion permutation and matched audits
in Online Supplement~C do not establish a robust residual.
\end{tablenotes}
\end{threeparttable}
\end{adjustbox}
\end{table}

\textit{Result.} The honest reading is deflationary. Exposure
alone already classifies cobidder contact well (exposure-only AUC
$0.905$ unconditionally; $0.713$ on the exposed subsample), and the
combined exposure-plus-score AUCs near $0.98$ overstate the score's
marginal value because $E_i$ encodes the label. Once exposure is
held fixed, the residual ordering is \emph{marginal at best and not
robust across designs}: the nested DeLong increment is only
$+0.010$ ($p=0.013$), the within-opportunity-stratum AUC for the
score is $0.471$ ($\approx$ chance), and the FL14 binary is no
better ($0.507$). The score does rank the binary excess-contact
target at AUC $0.764$ within MEDIUM cells, but the three companion
exercises in Online Supplement~C do not corroborate a robust
residual signal: a matched-stratum label permutation places the
observed PR-AUC at $p=0.127$ (not significant), a firm-exposure
simulation out-predicts the raw score on level PR-AUC ($p=1.000$),
the CEM opportunity-matched comparison gives a within-stratum AUC
of $0.626$ ($n=6{,}039$, full common support), and the within-stratum
change in $\Pr(\text{cobidder})$ for FL14 versus non-FL is in fact
\emph{negative} ($-0.017$). We therefore report no robust residual
ordering net of opportunity; the only marginally positive quantity
is the small nested increment, which the permutation tests do not
support as reliable.

\textit{Intensity-restricted label sensitivity.} A natural
objection is that the broad shared-tender-item label dilutes
genuine exposure with incidental single-contact cobidders, masking
a residual signal. We re-ran the full opportunity-adjusted battery
restricting positives to cobidders sharing at least two distinct
defendant tender-items ($368$ of the $651$ positives).
% src: outputs/sensitivity_contact2/ (02b_opportunity_sensitivity_contact2.R, 2026-06-04)
The restriction strengthens the deflationary conclusion rather
than reversing it: raw discrimination rises mechanically (ROC-AUC
$0.854$), but the exposure-only model rises faster ($0.873$ on the
exposed subsample; $0.956$ unconditionally) and again out-predicts
the score; the within-opportunity-stratum AUC remains at chance
($0.506$); the nested increment falls to $+0.003$ ($p=0.47$); and
the matched-stratum permutation again does not reject
($p=0.39$). The more tightly the label is anchored to intense
defendant contact, the more completely procurement opportunity
accounts for it.

\textit{Audit of the audit: leakage tiers, falsifiability, and
power.} Because a deflationary verdict can be manufactured as
easily as an inflated one, we audit the audit's own machinery.
% src: outputs/diagnostics/audit_armor/ (scripts 12/12b, 2026-06-04)
\emph{(i) Exposure-benchmark leakage tiers.} Contact-derived
benchmarks partially encode a contact-defined label. Ranking by
observed contact $O_i$ gives AUC $0.905$ and a plug-in $E_i$
$0.985$ (pure label encoding); the firm-leave-one-out $E_i$ used
in the main text gives $\valArmorExpLOO$; recomputing every
cell's defendant density from non-cobidder rows only---so no
eventual positive contributes to any firm's expected
contact---leaves a label-blind opportunity ranking of
$\valArmorExpLB$. The exposure benchmark is therefore an upper
bound on opportunity's role, and we rest no conclusion on
``exposure out-predicts the score'' alone.
\emph{(ii) Falsifiability of the within-stratum null.} Within the
same opportunity-decile strata that put the score at chance, a
positive-control score ($O_i$) ranks the label at
$\valArmorOiControlMedium$: the design detects a residual when
one exists. The within-stratum AUC for the score is stable across
cell granularities (coarse $\valArmorWithinCoarse$, medium
$\valArmorWithinMedium$, strict $\valArmorWithinStrict$; retained
support shrinks from $15{,}246$ to $615$ firms across that
range). Stratifying instead on the \emph{label-blind} expected
contact---strata that cannot mechanically absorb the
label---leaves the score at $\valArmorWithinLBall$; the negative
controls in Appendix~D identify that ordering as generic
co-participation geometry (placebo anchors reproduce it,
non-CADE high-volume-winner anchors exceed it), not
cartel-specific signal.
\emph{(iii) Power of the permutation design.} Injecting synthetic
within-stratum signal into the labels and re-running the
matched-strata permutation ($B=200$, $60$ simulations per effect
size) shows the test is correctly sized under the null (rejection
rate $0.05$ at zero injected signal) and rejects with probability
$\valArmorPowerTwo$ at a true within-stratum AUC of $0.52$,
$\valArmorPowerFive$ at $0.55$, and $\valArmorPowerTen$ at
$0.60$. The observed non-rejection therefore bounds any surviving
residual below a within-stratum AUC of roughly $0.55$.
\emph{(iv) Label-frozen timing.} Freezing the pool, the score,
and the label ingredients to $2009$--$2016$ (always-loser status
from training-window wins only; defendant contact from
training-window tender-items only) leaves $\valArmorFrozenPool$
rankable incumbents. The frozen score ranks new $2017$--$2019$
defendant contact at AUC $\valArmorFrozenProspAUC$
($\valArmorFrozenProspN$ positives) and fully frozen
retrospective contact at $\valArmorFrozenRetroAUC$
($\valArmorFrozenRetroN$ positives)---participation volume
forecasting generic future contact with active firms, consistent
with the anchor-placebo results and carrying no cartel-specific
content.

\subsection{Timing Information Sets and Strict Holdout}
\label{app:strict_holdout_submission}
\label{app:timing_sets_submission}

\textit{Purpose.} This block holds \emph{time} fixed: could a
regulator who froze the score \emph{before} the evaluation window
have ranked the firms later adjudicated as cobidders, without
peeking at future participation, wins, threshold, or CADE
adjudications?

\textit{Method.} The CADE files carry only judgment dates
(\texttt{data\_julgamento}, mostly 2015--2025) and \emph{no
conduct dates}, so all timing is \emph{year-level} (BEC
participation year), never day-level
(\texttt{DAY\_LEVEL\_TIMING\_UNAVAILABLE}). We define seven
candidate information sets, A--G, distinguished by which inputs
(score, zero-win status, threshold, label) are drawn pre- versus
post-window. We lead on the honest incumbent-triage designs --- C/D/G
(strict prospective score, rankable incumbent pool, retrospective
label, read as triage rather than prospective deployment) ---
report A as a labeled retrospective upper bound, B as a
leak-flagged decomposition, E to disclose the entrant coverage
hole, and document F (adjudication-observable-at-time) as
infeasible because all judgment dates post-date a 2009--2016
freeze. The strict holdout freezes the score, always-loser
status, and median$+1.5\times$IQR threshold on 2009--2016
(train-window threshold $\valStrictThreshTrain$, vs.\ full-window
$13.5$) and evaluates against 2017--2019 cobidder labels.
Table~\ref{tab:strict_holdout_submission} gives five nested
samples; the full A--G information-set register and the
leakage-channel cross-tabulation are in Online Supplement~C.

\begin{table}[htbp]
\centering
\caption{Strict Holdout $2009$--$2016 \to 2017$--$2019$ (Year-Level Timing)}
\label{tab:strict_holdout_submission}
\scriptsize
\begin{adjustbox}{max width=\textwidth,center}
\begin{threeparttable}
\begin{tabular}{lrrrrrrrr}
\toprule
Sample & $N$ & Pos. & ROC AUC & PR-AUC & Prec@500 & Recall@500 & Lift@500 & Tie-at-0 \\
\midrule
1. Full test candidate & $41{,}444$ & $651$ & $0.474$ & $0.014$ & $0.000$ & $0.000$ & $0.00$ & $0.246$ \\
2. Rankable ($T^{\text{train}}{>}0$) & $32{,}682$ & $498$ & $0.471$ & $0.013$ & $0.000$ & $0.000$ & $0.00$ & $0.043$ \\
3. Training always-loser & $21{,}819$ & $651$ & $0.684$ & $0.077$ & $0.124$ & $0.095$ & $4.16$ & $0.402$ \\
4. Training AL \& test-active & $11{,}351$ & $312$ & $0.667$ & $0.084$ & $0.142$ & $0.228$ & $5.17$ & $0.772$ \\
5. Zero-win both (\textsc{leaks}) & $16{,}843$ & $651$ & $0.646$ & $0.086$ & $0.136$ & $0.104$ & $3.52$ & $0.280$ \\
\bottomrule
\end{tabular}
\begin{tablenotes}
\scriptsize
\item \textit{Notes:} Continuous loser-side score
($\log T^{\text{train}}$). Sample 3 is the honest strict pool;
within it, the FL-binary AUC is $\valStrictFLAUC$ and the
continuous AUC is $\valStrictContAUC$ (threshold
$=\valStrictThreshTrain$, vs.\ full-window $13.5$). Sample 5 leaks
future wins via full-period always-loser status and is reported as
a leak-flagged sensitivity only. \textbf{On the full candidate
universe (sample 1) the strict screen collapses to chance:} ROC
$\valStrictFullAUC$ (below chance), PR-AUC $0.014$, and
precision@$500=$recall@$500=\valStrictFullPrec$, with about $25\%$ of
firms tied at score zero. Of the $\valStrictRankablePos$ rankable
positives and $\valMainEntrantPositives$ entrant positives
($\valEntrantPosShare$ of $\valMainCobidders$), entrants form a
structural blind spot the prospective score cannot reach. ``Tie-at-0''
is the share of the candidate set carrying score zero. Source:
\texttt{table\_D\_strict\_2009\_2016\_to\_2017\_2019.csv};
\texttt{strict\_holdout\_composition.csv}.
\end{tablenotes}
\end{threeparttable}
\end{adjustbox}
\end{table}

\textit{Result.} The reading is deliberately deflationary. Only a
limited within-pool residual ordering remains (sample 3, AUC
$\valStrictContAUC$--$\valStrictFLAUC$), but on the full
firm universe the strict screen has no operational traction
(precision@$500=\valStrictFullPrec$): the entrant coverage hole
and the mass of score-zero ties are the proximate causes. Two
companion exercises in Online Supplement~C confirm the pattern. A
rolling-origin validation expanding the training window for
$t=2014,\dots,2019$ posts PR-AUC $\approx 0.012$--$0.014$ and
precision@$500=$recall@$500=0$ on the full universe in every origin
year, with the worst ROC below chance ($\valRollWorstAUC$, 2015);
the within-pool (training-always-loser) ROC stays in the
$0.66$--$0.70$ range across origins, again an artifact of a thicker
retrospective training history rather than real-time learning. A leakage decomposition of the in-sample item-level
reference AUC ($\valLeakRefAUC$) confirms that the surviving
structural component stays in the
$\valLeakStructLow$--$\valLeakStructHigh$ range under
cobidder-firm holdout (out-of-fold AUC $\valLeakCVAUC$) and
temporal holdout (AUC $\valAUCFLfirmTemp$), with the
$\valLeakTautLow$--$\valLeakTautHigh$ AUC lost to mechanical reuse
disclosed as such.

\subsection{Case Composition}
\label{app:loco_submission}
\label{app:environment_dominance_submission}

\textit{Purpose.} A generalization concern is whether the pooled
signal is carried by one or two adjudicated cases rather than the
construct.

\textit{Method.} We link six cases with enough cobidders to score
(a seventh, \texttt{sacos\_de\_lixo}, is too sparse and is
flagged) and re-evaluate the ranking after dropping the largest
and the top-two cases, and under case-balanced weighting.
Table~\ref{tab:case_dominance_submission} reports the
leave-largest-case-out collapse.

\begin{table}[htbp]
\centering
\caption{Case Dominance: Leave-Largest-Case-Out}
\label{tab:case_dominance_submission}
\scriptsize
\begin{threeparttable}
\begin{tabular}{lrrrrr}
\toprule
Scenario & Pos. & ROC AUC & PR-AUC & Prec@500 & Rec@500 \\
\midrule
Full (pooled) & $651$ & $0.761$ & $\valLOCOfullPRAUC$ & $\valFullPrec$ & $0.166$ \\
Drop largest case & $443$ & $0.763$ & $\valLOCOdropLargestPRAUC$ & $\valLOCOdropLargestPrec$ & $0.149$ \\
Drop top-two cases & $268$ & $0.752$ & $0.058$ & $0.092$ & $0.172$ \\
Case-balanced (mean over cases) & --- & --- & --- & $\valCaseBalancedPrec$ & $0.173$ \\
\bottomrule
\end{tabular}
\begin{tablenotes}
\scriptsize
\item \textit{Notes:} Dropping the single largest case
(\texttt{trens\_metros}) lowers PR-AUC from
$\valLOCOfullPRAUC$ to $\valLOCOdropLargestPRAUC$ ($-37\%$) and
precision@$500$ from $\valFullPrec$ to $\valLOCOdropLargestPrec$;
dropping the top two cases takes PR-AUC to $0.058$.
\textbf{Case-balanced precision@$500$ (weighting each case equally
rather than each firm) is $\valCaseBalancedPrec$, against the
firm-weighted pooled $\valFullPrec$:} the pooled operational number
is concentrated in the largest case. ROC barely moves (it is
uninformative at this prevalence). Source:
\texttt{table\_H\_case\_dominance\_validation.csv}.
\end{tablenotes}
\end{threeparttable}
\end{table}

\textit{Result.} Dropping the single largest case lowers PR-AUC
from $\valLOCOfullPRAUC$ to $\valLOCOdropLargestPRAUC$ ($-37\%$);
case-balanced precision@$500$ is $\valCaseBalancedPrec$ against the
firm-weighted pooled $\valFullPrec$, so the pooled operational
number is concentrated in the largest case (which holds
$\valTopCasePosShare$ of positives and $\valTopCaseTP$ of the
true positives in the top $500$). The concentration is structural
across the environment, documented in Online Supplement~C: the
largest item group holds $\valTopItemGroupShare$ of positives
(item-group HHI $0.100$) while the buyer distribution is diffuse
(buyer HHI $0.009$), and dropping the largest item group leaves
PR-AUC essentially unchanged ($0.143 \to 0.134$). A clustered
randomization-inference test permuting cobidder labels within
item-group $\times$ year clusters ($B=1{,}000$) clears the null on
the ranking metrics (ROC, PR-AUC, precision@$500$ all
$p=\valClusterRIp$), but top-$500$ case coverage is \emph{not}
significant ($p=\valClusterRIcovP$), consistent with one case
anchoring the operational signal.

\subsection{Direct-Defendant Scope Check}
\label{app:direct_scope_submission}
\label{app:direct_defendant_timing_submission}

If the statistic were a generic cartel-membership detector, it
should classify direct CADE defendants as well as cobidders. It
does not: the FL-binary score (retrospective incumbent triage)
ranks direct defendants at ROC $\approx \valDirectStrictAUC$ under
both full and strict timing ($0.491$ and $0.492$), i.e.\ at chance,
with precision@$500=0$. Continuous participation volume does carry
some information about defendants---ROC $0.66$--$0.70$, modestly
above chance---but with precision@$500\le0.004$, so it does not
operationally retrieve them either (full grid in Online
Supplement~C). The asymmetry is the design: the loser-side
\emph{binary} flag is uninformative about defendants, who typically
sit on the winner side of the arrangement, while participation
volume is not entirely silent. This is a scope boundary, not an
escape hatch.

\subsection{Limitations}
\label{app:opp_limitations_submission}
\label{app:timing_limitations_submission}

Several limitations bound this appendix. The opportunity cells are
\emph{award-layer proxies} built from item, year, and
procuring-unit identifiers, not a full description of each firm's
competitive opportunity set, and procurement modality is dropped
because it is absent from that layer. All timing is
\emph{year-level} (\texttt{DAY\_LEVEL\_TIMING\_UNAVAILABLE}): the
CADE files carry only judgment dates (2015--2025), no conduct
dates, so we cannot align the screen to when a cartel operated and
the adjudication-observable design is infeasible rather than
estimated. There is an \emph{entrant coverage hole}
($\valMainEntrantPositives/\valMainCobidders$,
$\valEntrantPosShare$, of positives carry score zero with no
training history) and the signal is \emph{concentrated} by case
(dropping the largest cuts PR-AUC by $37\%$), item group, and year,
so the pooled operational number overstates breadth. ROC is a red
herring at sub-$1\%$ prevalence: precision and recall---not
ROC---are the operational metrics, and they stay low under the
pooled diagnostic and collapse to zero on the full strict-timing
universe. Most importantly, the
CADE-defendant labels are \emph{adjudication-anchored exposure}
labels, not membership findings: a tagged firm shared a
tender-item with an adjudicated defendant; it is not thereby a
cartel member, a party to any agreement, or liable for anything.
Net of opportunity and timing, the residual ordering is
\emph{marginal at best and not robust across designs}; what the
award layer reliably delivers is a coarse, retrospective basis for
prioritizing scrutiny, and it is \textbf{not legal proof, not
identification of cartel members, and not an estimate of damages
or causal effect}.
